% Paper A: introduction (follows the author's outline).
\section{Introduction}\label{sec:intro}
A partial differential equation with constant coefficients becomes, in Fourier variables, a polynomial: the operator
$(\partial_t+c\,\partial_x)(\partial_t-\nu\,\partial_x^2)$ of a quantity that is carried by a flow and diffuses becomes
$(\tau+c\,\xi)(\tau-\nu\,\xi^2)$. Polynomials can be multiplied and factored, and a factorization of the symbol splits
every solution into pieces, one for each factor: here a travelling profile, whose Fourier modes lie on the line
$\lambda=-ic\kappa$, and a diffusing one, whose modes lie on the parabola $\lambda=-\nu\kappa^2$ (\cref{fig:branches}).
This paper is about that splitting, which we call the \emph{grouping structure} of the equation, and about how a
machine-learning algorithm can find it in noisy data and use it to recover the equation. The example matters in
practice: methods that regress $u_t$ on a library of terms look for a single sum of terms, and every one we tested
returns a wrong equation for it at $5\%$ noise, without warning.

% Shared figure for the introductions: the two groups of the running example as curves of modes.
\begin{figure}[tbp]
\centering
\begin{tikzpicture}[x=0.40cm,y=0.40cm,>=Stealth,font=\footnotesize]
  % (a) oscillation slice
  \begin{scope}
    \draw[->,gray] (-4.5,0)--(4.8,0) node[right,black]{$\kappa$};
    \draw[->,gray] (0,-4.5)--(0,4.8) node[above,black]{$\operatorname{Im}\lambda$};
    \draw[very thick,color={rgb,255:red,42;green,120;blue,214}] (-4.2,4.2)--(4.2,-4.2);
    \draw[very thick,color={rgb,255:red,235;green,104;blue,52},dashed] (-4.2,0)--(4.2,0);
    \foreach \k in {-3,-2,-1,1,2,3} \fill[color={rgb,255:red,42;green,120;blue,214}] (\k,-\k) circle (2.6pt);
    \foreach \k in {-2,-1,1,2} \fill[color={rgb,255:red,235;green,104;blue,52}] (\k,0) +(-2.4pt,-2.4pt) rectangle +(2.4pt,2.4pt);
    \node[anchor=west,align=left] at (1.2,3.3) {transport\\$\lambda=-ic\kappa$};
    \node at (0,-5.8) {(a) oscillation};
  \end{scope}
  % (b) damping slice
  \begin{scope}[xshift=10.6cm]
    \draw[->,gray] (-4.5,0)--(4.8,0) node[right,black]{$\kappa$};
    \draw[->,gray] (0,-4.5)--(0,1.6) node[above,black]{$\operatorname{Re}\lambda$};
    \draw[very thick,color={rgb,255:red,42;green,120;blue,214},dashed] (-4.2,0)--(4.2,0);
    \draw[very thick,color={rgb,255:red,235;green,104;blue,52},domain=-3.8:3.8,samples=60] plot (\x,{-0.3*\x*\x});
    \foreach \k in {-3,-2,-1,1,2,3} \fill[color={rgb,255:red,42;green,120;blue,214}] (\k,0) circle (2.6pt);
    \foreach \k in {-2,-1,1,2} \fill[color={rgb,255:red,235;green,104;blue,52}] (\k,{-0.3*\k*\k}) +(-2.4pt,-2.4pt) rectangle +(2.4pt,2.4pt);
    \node[anchor=west,align=left] at (1.8,1.6) {diffusion\\$\lambda=-\nu\kappa^2$};
    \node at (0,-5.8) {(b) damping};
  \end{scope}
\end{tikzpicture}
\caption{The idea of the paper on the equation $(\partial_t+c\,\partial_x)(\partial_t-\nu\,\partial_x^2)u=0$. Every Fourier mode
$e^{\lambda t+i\kappa x}$ of a solution lies on the zero set of one factor of the symbol $(\tau+c\xi)(\tau-\nu\xi^2)$: circles
on the line of the transport factor, squares on the parabola of the diffusion factor (schematic; \cref{fig:directions}
shows the modes found in noisy data). Learning the equation means learning these two groups.}\label{fig:branches}
\end{figure}


Two questions organize the work. The first is analytical: which pieces of an equation does a solution actually carry,
how uniquely, and how stably? The answer is that the pieces are the \emph{factors} of the symbol, not the additive
\emph{terms} of the operator; that the split is unique up to a finite-dimensional overlap whose dimension is an
intersection number; and that its stability is governed by the Friedrichs angle between the solution spaces of the
factors. The second question is algorithmic: how can one learn the groups from finitely many noisy observations, and
what should the output say about how well they were learned? Our answer is an algorithm that groups twice. Within one
sample of the field it groups Fourier modes into factors, using the fact that the modes of one factor lie on one curve
of the dispersion relation. Across posterior samples it groups the factors found in each sample into clusters, using a
Dirichlet-process mixture. The output is a distribution over equations with labels that say what the data can and
cannot support.

\paragraph{Contributions}
\begin{enumerate}
  \item An analytical account of the grouping structure: identification and uniqueness (\cref{thm:ident,prop:unique}),
  the factor decomposition with its exact ambiguity $\deltaP$ (\cref{thm:factor}), its conditioning
  (\cref{thm:hm}), and the impossibility of grouping by terms (\cref{prop:terms,prop:discont}).
  \item An algorithm that turns this structure into a learning procedure: interpolation with a sampling-based
  treatment of the high-frequency band, a matrix pencil per wavenumber, a joint grouping-and-fitting of the modes, and
  a Bayesian layer that groups factors across samples and attaches intervals whose empirical coverage is
  $0.94$ at nominal $0.90$.
  \item Evidence: on twenty random equations not used in development the method is right or equivalent on $19$
  ($14$ with a confident label, $5$ with an abstaining label and correct factors) and never wrong without warning,
  against $13$ or more such failures for each of the three library-regression baselines run on this set; on $60$ grid
  inputs with independent fields and noise it is always right, against $13$ silent failures for weak SINDy; on a
  benchmark of ten equations in three designs, on which its parameters were tuned, it recovers all $30$ factor
  sets, while the best library regression is right or equivalent on $25$ and wrong without warning on $5$.
\end{enumerate}

% Shared: relation to prior work (what is classical, what is new).
\paragraph{Relation to prior work}
Most methods that learn a PDE from data regress a time derivative on a library of candidate terms: SINDy and
PDE-FIND~\cite{BruProKut16,RudBruProKut17}, their weak, ensemble and Bayesian variants~\cite{MesBor21,FasKutBru22,HirBarKut22},
and Gaussian-process formulations of linear operators~\cite{RaiPerKar17}. Their unit is the \emph{term}. Dynamic mode
decomposition~\cite{Sch10,SasKut22} works with eigenvalues, as we do, but returns a linear operator on the data, not an
equation, and does not group eigenvalues by the factor they belong to. Exponential fitting by matrix
pencils~\cite{HuaSar90} is a classical tool of signal processing; the algebra of constant-coefficient operators, the
Ehrenpreis--Palamodov fundamental principle and the solvability results behind \cref{thm:factor} are classical
analysis~\cite{Ehr70,Pal70,Mal56,Hor90,Obe90}; Dirichlet-process mixtures are standard Bayesian
clustering~\cite{BleJor06}. What is new here is the connection between them. (i) We identify the \emph{factor} of the
symbol, not the term, as the unit a learning method should group by, and we prove that, for two terms whose varieties meet in finitely many points, per-term grouping is impossible
outside a finite-dimensional set of solutions (\cref{prop:terms,prop:discont}). (ii) We state what finite data can and
cannot certify about such groups (\cref{thm:ident,prop:certify}) and give the exact conditioning of the
separation between groups (\cref{thm:hm,prop:bayes-angle}). (iii) We build an algorithm that groups twice, modes into
factors within a sample and factors across posterior samples, and whose output carries these limits as labels and
intervals.


\paragraph{Scope} Everything proved and implemented concerns scalar linear equations with constant coefficients in one
time and one space variable, on rectangular windows. \Cref{sec:nonlinear} explains which parts of the structure are lost
for nonlinear equations, and \cref{rem:varcoef} how the spectral step reads variable coefficients; neither is claimed as a
result.

\paragraph{Organization} \Cref{sec:analysis} develops the analytical framework: the operator ring, exponential-polynomials,
identification, the factor decomposition and its spectral projectors, the Friedrichs angle, and the interaction term with
its energy. \Cref{sec:nonlinear} discusses nonlinear equations. \Cref{sec:algorithm} describes the deterministic
algorithm and \cref{sec:bayesian} its statistical layer. \Cref{sec:results} follows one data set through every stage and
reports benchmarks. \Cref{sec:limits} separates the inference problems that have a remedy from those that do not, and
\cref{sec:conclusion} concludes. \Cref{fig:roadmap} summarizes the method; proofs are in \cref{app:proofs}. \emph{Reader's guide.} The algorithm uses four results: the factor decomposition (\cref{thm:factor}), the
stability of the split (\cref{thm:hm}), the counting condition (\cref{prop:counting}) and the limits of finite data
(\cref{prop:certify}), which explain the labels. A reader interested mainly in the method can skim the material on
exponential-polynomials and identification on a first reading.

% Shared figure: roadmap of the method, each box with the section and the result it rests on.
\begin{figure}[tbp]
\centering
\begin{tikzpicture}[
  node distance=3.5mm and 4.5mm,
  box/.style={draw=black!55, rounded corners=2pt, align=center, font=\footnotesize, inner sep=3pt,
              minimum height=10mm, text width=20mm, fill=white},
  grp/.style={box, fill=blue!6},
  thy/.style={font=\scriptsize\itshape, text=black!60, align=center},
  arr/.style={-{Stealth[length=2mm]}, black!60, thick}]
\node[box] (data) {noisy data\\$\{(t_i,x_i,y_i)\}$};
\node[box, right=of data] (gp) {field samples\\$u^{(1)},\dots,u^{(\Nsamp)}$};
\node[box, right=of gp] (modes) {modes $(\kappa,\lambda)$\\per wavenumber};
\node[grp, right=of modes] (fac) {factors $Q_j$\\(groups of modes)};
\node[grp, below=8mm of fac] (cloud) {cloud of factors\\over samples};
\node[box, left=of cloud] (clus) {clusters and\\$P(k)$, $P(\Mfit)$};
\node[box, left=of clus] (split) {components $u_j$,\\$w_j$, $\gN$, $\rho$};
\node[box, left=of split] (out) {equation with\\intervals and label};
\draw[arr] (data)--(gp); \draw[arr] (gp)--(modes); \draw[arr] (modes)--(fac);
\draw[arr] (fac)--(cloud); \draw[arr] (cloud)--(clus); \draw[arr] (clus)--(split); \draw[arr] (split)--(out);
\node[thy, above=1mm of gp] {\S\ref{sec:alg-interp}--\ref{sec:nyquist}};
\node[thy, above=1mm of modes] {\S\ref{sec:alg-pencil}, Lem.~\ref{lem:ep}};
\node[thy, above=1mm of fac] {\S\ref{sec:alg-structfit}, Thm.~\ref{thm:factor}};
\node[thy, below=1mm of cloud] {\S\ref{sec:bayes-cloud}};
\node[thy, below=1mm of clus] {\S\ref{sec:alg-loop}, \S\ref{sec:bayes-cloud}};
\node[thy, below=1mm of split] {Thms.~\ref{thm:hm}, \ref{thm:eint}};
\node[thy, below=1mm of out] {\S\ref{sec:alg-labels}, Thm.~\ref{thm:ident}};
\end{tikzpicture}
\caption{The method at a glance. Shaded: the two grouping steps, of modes into factors within one sample and of
factors across samples. Each box names the section that describes it and the result that justifies it.}\label{fig:roadmap}
\end{figure}

